% !PRINTED 385
\weberpage{385}
%<<CONT
divides. In this case, however, $(D, p^s) = (\Delta, p^s)$ [\S\,85 (9)],
and so what has to be proved is:
\[
\chi(\delta, p^k) \sum_{0,k}^{s} (\Delta, p^s) = \sum_{0,k}^{s} (\delta, p^s)(\delta', p^{k-s}).
\]
If now $p$ does not divide $\delta$, then
\[
\chi(\delta, p^k) = (\delta, p^k),
\]
and if on the right one multiplies each term under the summation sign by
$(\delta, p^{k-s})^2 = 1$, one obtains, in agreement with
the left-hand side,
\[
(\delta, p^k) \overset{s}{\sum} (\Delta, p^{k-s}). % [sic: lower limit omitted in source]
\]
But if $p$ divides $\delta$, and consequently does not divide $\delta'$, then $\chi(\delta, p^k)$
$= (\delta', p^k)$; and if one multiplies in the same way by $(\delta', p^s)^2$, the same
result follows, whereby (22) is proved.

8. \gesperrt{If $\chi(\delta, n)$ is different from $\chi(\delta', n)$, then}
\weq{23}{\Sigma(D, \bar e) = 0, \quad \Sigma(\delta, e)(\delta', e') = 0.}
For since the two sums do not change when $\delta$ is interchanged with
$\delta'$, this follows from (22).

9. \gesperrt{If $\delta_1, \delta_2$ are two stem divisors of $D$ and $\delta_1\delta_2$ their
symbolic product, then}
\weq{24}{\chi(\delta_1, n)\, \chi(\delta_2, n) = \chi(\delta_1\delta_2, n).}
Because of (21) this formula needs to be proved only for the case
that $n$ is a prime.

But symbolically we have
\[
(\delta_1\delta_2)' = \delta_1'\delta_2 = \delta_1\delta_2',
\]
and from this (24) follows easily from the definition (19).

\begin{center}
§ 105.\quad \gesperrt{Application of the Legendre Symbol.}
\end{center}

If one wishes to use the more familiar notation of \gesperrt{Legendre}
and \gesperrt{Jacobi}, the matter takes the following shape.
If $l$ denotes an odd prime and $a$ a number not
divisible by $l$, then according to \gesperrt{Legendre}:
\weq{1}{\left(\frac{a}{l}\right) = +1,}
\gesperrt{when $a$ is a quadratic residue of $l$,}
\[
\left(\frac{a}{l}\right) = -1,
\]
%CONT>>
